\documentclass[10pt,a4paper]{article}

% ==============================================================================
% TEMPLATE EDITORIAL — PERIÓDICO EM 2 COLUNAS (DOSSIÊ ANUAL E CARTA DE CONJUNTURA)
% JORNAL ECONÔMICO DO TOCANTINS
%   - Dossiê Anual (DA): Verde-Floresta (#0B5D32), Marfim (#FAF6EE) e Carmim (#A61C1C)
%   - Carta de Conjuntura (CC): Azul Cobalto (#0A2E73), Esmeralda (#1B7A43),
%     Lima Editorial (#8CE000 no selo escuro) e Azul-Gelo (#F3F6FB)
% ==============================================================================
\usepackage[utf8]{inputenc}
\usepackage[T1]{fontenc}
\usepackage[brazil]{babel}
\usepackage{helvet}
\renewcommand{\familydefault}{\sfdefault}

\usepackage[top=1.35cm, bottom=1.45cm, left=1.35cm, right=1.35cm, headheight=14pt, headsep=8pt, footskip=18pt]{geometry}
\usepackage{xcolor}
\usepackage{graphicx}
\usepackage{booktabs}
\usepackage{tabularx}
\usepackage{longtable}
\usepackage{array}
\usepackage{colortbl}
\usepackage{tikz}
\usepackage{fancyhdr}
\usepackage{lastpage}
\usepackage{enumitem}
\usepackage{microtype}
\usepackage{setspace}
\usepackage{multicol}
\usepackage{amssymb}
\usepackage{hyperref}

% Paleta Editorial Dinâmica (Proporção 75% Neutro / 20% Primária / 5% Acento)
\definecolor{corprimaria}{HTML}{0F2537}
\definecolor{corsecundaria}{HTML}{334155}
\definecolor{coracento}{HTML}{9A6C11}
\definecolor{coracentoheader}{HTML}{C8942B}
\definecolor{coracentotexto}{HTML}{0F2537}
\definecolor{corfundo}{HTML}{FAF8F5}
\definecolor{cinzatexto}{HTML}{1E293B}
\definecolor{cinzasub}{HTML}{475569}
\definecolor{cinzaborda}{HTML}{CBD5E1}
\definecolor{verdealta}{HTML}{15803D}
\definecolor{vermelhobaixa}{HTML}{B91C1C}

\hypersetup{
  colorlinks=true,
  linkcolor=corprimaria,
  urlcolor=corsecundaria,
  pdftitle={Dossiê Anual — Arrecadação Tributária e Finanças Públicas do Tocantins (2023–2026)},
  pdfauthor={Jornal Economico do Tocantins}
}

\setlength{\parindent}{0pt}
\setlength{\parskip}{3.5pt}
\setlength{\columnsep}{0.60cm}
\setlength{\columnseprule}{0.4pt}
\def\columnseprulecolor{\color{cinzaborda}}
\renewcommand{\arraystretch}{1.14}

% Cabeçalho e Rodapé de Periódico Econômico (Sem sobreposição central)
\pagestyle{fancy}
\fancyhf{}
\renewcommand{\headrulewidth}{0pt}
\renewcommand{\footrulewidth}{0.5pt}
\renewcommand{\footrule}{\hbox to\headwidth{\color{coracento}\leaders\hrule height \footrulewidth\hfill}}

\fancyhead[L]{\ifnum\value{page}>1 \scriptsize\bfseries\color{corprimaria}JORNAL ECONÔMICO DO TOCANTINS \enspace{\color{coracento}$\vert$}\enspace {\mdseries\color{cinzasub}DOSSIÊ ANUAL / SEMESTRAL (DA)}\fi}
\fancyhead[R]{\ifnum\value{page}>1 \scriptsize\color{cinzasub}Ref.: \textbf{Acumulado no Exercício e Série Plurianual (2023–2026)} $\cdot$ \#07/2026\fi}

\fancyfoot[L]{\scriptsize\color{cinzasub}\textbf{\color{corprimaria}Jornal Econômico do Tocantins} $\cdot$ Lorenzo Costa | Graduado em Economia — Jornal Econômico do Tocantins $\cdot$ Publicado em 14/11/2026}
\fancyfoot[R]{\scriptsize\bfseries\color{corprimaria}Pág. \thepage\ de \pageref*{LastPage}}

% Banner Editorial Limpo de Abertura de Seção Principal (Caderno Periódico)
\newcommand{\bannersecaoperiodico}[3]{%
  \par\vspace{0.16cm}%
  \noindent{\color{corprimaria}\hrule height 1.4pt}%
  \par\vspace{4pt}%
  \noindent{\scriptsize\bfseries\color{coracento}SEÇÃO #1 \enspace$\cdot$\enspace}{\normalsize\bfseries\color{corprimaria}#2}%
  \par\vspace{2pt}%
  {\scriptsize\itshape\color{cinzasub}#3\par}%
  \vspace{3pt}%
  {\color{cinzaborda}\hrule height 0.5pt}%
  \par\vspace{0.10cm}%
}

% Subtítulo de Coluna dentro do Layout de 2 Colunas
\newcommand{\secaoperiodico}[1]{%
  \par\vspace{0.10cm}%
  \noindent{\small\bfseries\color{corprimaria}#1}%
  \par\vspace{0.03cm}%
  {\color{coracento}\hrule height 0.6pt}%
  \vspace{0.06cm}%
}

\begin{document}
\color{cinzatexto}

% ==============================================================================
% MASTHEAD EDITORIAL CLARO DE PERIÓDICO (DOSSIÊ ANUAL OU CARTA DE CONJUNTURA)
% ==============================================================================
\noindent{\color{corprimaria}\hrule height 2.2pt}%
\par\vspace{2pt}%
{\color{coracento}\hrule height 0.8pt}%
\par\vspace{6pt}%
\noindent\hbox to \linewidth{%
  {\scriptsize\bfseries\color{coracento}\MakeUppercase{DOSSIÊ ANUAL / SEMESTRAL (DA)} \enspace$\cdot$\enspace \#07/2026}%
  \hfill
  {\scriptsize\color{cinzasub}\textbf{\color{corprimaria}JORNAL ECONÔMICO DO TOCANTINS} \enspace$\vert$\enspace Publicado em 14/11/2026}%
}%
\par\vspace{5pt}%
{\fontsize{14.5}{17.4}\selectfont\bfseries\color{corprimaria}Dossiê Anual — Arrecadação Tributária e Finanças Públicas do Tocantins (2023–2026)\par}%
\vspace{3pt}%
{\footnotesize\itshape\color{cinzasub}Análise acumulada de 12 meses e plurianual das receitas próprias (ICMS, IPVA e outros tributos) e das transferências constitucionais do FPE \enspace$\vert$\enspace \textbf{\color{corprimaria}Ref.: Acumulado no Exercício e Série Plurianual (2023–2026)}\par}%
\par\vspace{6pt}%
{\color{cinzaborda}\hrule height 0.6pt}%

\vspace{0.12cm}

% ==============================================================================
% CORPO COMPLETO DO PERIÓDICO (DOSSIÊ ANUAL OU CARTA DE CONJUNTURA)
% ==============================================================================
\noindent\makebox[\textwidth][s]{%
\begin{tikzpicture}[baseline=(box.base)]
  \node[fill=white, draw=cinzaborda, line width=0.35pt,
        inner xsep=6pt, inner ysep=4pt, text width=0.206\textwidth, minimum height=1.14cm, align=left] (box) {%
    {\fontsize{6.0}{7.2}\selectfont\bfseries\color{cinzasub}RECEITA TOTAL DO MÊS}\par\vspace{1.2pt}%
    {\fontsize{9.2}{10.8}\selectfont\bfseries\color{corprimaria}R\$ 1,19 Bi}\par\vspace{1pt}%
    {\fontsize{5.8}{6.8}\selectfont\color{cinzasub}Ref.: 09/26}%
  };
  \fill[coracento] (box.north west) rectangle ([yshift=-1.6pt]box.north east);
\end{tikzpicture}\hfill%
\begin{tikzpicture}[baseline=(box.base)]
  \node[fill=white, draw=cinzaborda, line width=0.35pt,
        inner xsep=6pt, inner ysep=4pt, text width=0.206\textwidth, minimum height=1.14cm, align=left] (box) {%
    {\fontsize{6.0}{7.2}\selectfont\bfseries\color{cinzasub}ARRECADAÇÃO DE ICMS}\par\vspace{1.2pt}%
    {\fontsize{9.2}{10.8}\selectfont\bfseries\color{corprimaria}R\$ 566,48 Mi}\par\vspace{1pt}%
    {\fontsize{5.8}{6.8}\selectfont\color{cinzasub}Var. Mês: -5,25\% | YoY: +1.1\%}%
  };
  \fill[coracento] (box.north west) rectangle ([yshift=-1.6pt]box.north east);
\end{tikzpicture}\hfill%
\begin{tikzpicture}[baseline=(box.base)]
  \node[fill=white, draw=cinzaborda, line width=0.35pt,
        inner xsep=6pt, inner ysep=4pt, text width=0.206\textwidth, minimum height=1.14cm, align=left] (box) {%
    {\fontsize{6.0}{7.2}\selectfont\bfseries\color{cinzasub}REPASSES DO FPE}\par\vspace{1.2pt}%
    {\fontsize{9.2}{10.8}\selectfont\bfseries\color{corprimaria}R\$ 438,14 Mi}\par\vspace{1pt}%
    {\fontsize{5.8}{6.8}\selectfont\color{cinzasub}Tesouro Nacional / STN}%
  };
  \fill[coracento] (box.north west) rectangle ([yshift=-1.6pt]box.north east);
\end{tikzpicture}\hfill%
\begin{tikzpicture}[baseline=(box.base)]
  \node[fill=white, draw=cinzaborda, line width=0.35pt,
        inner xsep=6pt, inner ysep=4pt, text width=0.206\textwidth, minimum height=1.14cm, align=left] (box) {%
    {\fontsize{6.0}{7.2}\selectfont\bfseries\color{cinzasub}IPVA E OUTROS TRIBUTOS}\par\vspace{1.2pt}%
    {\fontsize{9.2}{10.8}\selectfont\bfseries\color{corprimaria}R\$ 189,35 Mi}\par\vspace{1pt}%
    {\fontsize{5.8}{6.8}\selectfont\color{cinzasub}IPVA puro: R\$ 53,89 Mi}%
  };
  \fill[coracento] (box.north west) rectangle ([yshift=-1.6pt]box.north east);
\end{tikzpicture}%
}

\vspace{0.03cm}
\par\vspace{0.05cm}
\noindent{\color{cinzaborda}\hrule height 0.5pt}\par\vspace{4pt}
{\fontsize{7.8}{9.4}\selectfont\bfseries\color{corprimaria}\MakeUppercase{Síntese dos Indicadores — Balanço em 12 Meses}}\par\vspace{2pt}
{\scriptsize\color{cinzatexto}
\begin{itemize}[leftmargin=10pt, itemsep=1.2pt, topsep=1pt, parsep=0pt]
   \item \textbf{\color{corprimaria}Arrecadação de ICMS (09/26):} Arrecadação de ICMS (09/26) sofreu retração acentuada de -5,25\% (queda drástica que requer atenção conjuntural). (patamar atual: R\$ 566,48 Mi)
   \item \textbf{\color{corprimaria}Cota-Parte do FPE (Transferência União):} Cota-Parte do FPE (Transferência União) sofreu retração acentuada de -19,32\% (queda drástica que requer atenção conjuntural). (patamar atual: R\$ 438,14 Mi)
   \item \textbf{\color{corprimaria}IPVA e Outros Tributos Estaduais:} IPVA e Outros Tributos Estaduais sofreu retração acentuada de -10,32\% (queda drástica que requer atenção conjuntural). (patamar atual: R\$ 189,35 Mi)
   \item \textbf{\color{corprimaria}Variação Interanual do ICMS (YoY):} Variação Interanual do ICMS (YoY) avançou +1,14\% em trajetória de expansão gradual. (patamar atual: +1.1\%)
\end{itemize}
}
\noindent{\color{cinzaborda}\hrule height 0.5pt}\par\vspace{0.05cm}

\vspace{0.03cm}
\begin{multicols}{2}
\secaoperiodico{1. Análise Estrutural Anual / Acumulado em 12 Meses}
{\small Sob a ótica anual e plurianual (2023–2026), o desempenho fiscal do Estado do Tocantins demonstra fortalecimento gradual da arrecadação tributária própria frente às transferências federais. Enquanto os repasses constitucionais do Fundo de Participação dos Estados (FPE) refletem a arrecadação nacional de IR e IPI, o ICMS estadual mantém crescimento acumulado positivo de 3,65\% em 2026, superando o patamar observado nos exercícios anteriores.

\medskip
A composição setorial do ICMS evidencia que Combustíveis, Comércio e Indústria respondem por mais de 78\% da arrecadação própria estadual, funcionando como termômetro direto da circulação de mercadorias no corredor logístico Norte-Sul e do consumo das famílias nos municípios polos.}
\end{multicols}
\begin{center}\includegraphics[width=0.84\textwidth]{figuras/2026_11_dossie_financas_fin_receitas.pdf}\end{center}
\par\vspace{0.08cm}
\noindent\begin{minipage}[t]{0.482\textwidth}
\par\vspace{0.08cm}
\noindent\begin{minipage}{\linewidth}
  \noindent{\fontsize{7.8}{9.5}\selectfont\bfseries\color{corprimaria}\rule{0.10cm}{0.20cm}\hspace{4pt}Execução Mensal das Receitas Próprias e Transferências Constitucionais do Tocantins}\par\vspace{2pt}
  {\scriptsize\setlength{\tabcolsep}{4.0pt}%
\noindent\resizebox{\linewidth}{!}{%
\begin{tabular}{lrrrr}
\toprule
\rowcolor{corprimaria} \textbf{\color{white}Mês (2026)} & \textbf{\color{white}ICMS (SEFAZ-TO)} & \textbf{\color{white}Cota-Parte FPE (STN)} & \textbf{\color{white}IPVA + Outros Tributos} & \textbf{\color{white}Receita Total Mensal} \\
\midrule
Jan & R\$ 531,60 Mi & R\$ 608,22 Mi & R\$ 199,28 Mi & R\$ 1.339,10 Mi \\
\rowcolor{corfundo} Fev & R\$ 465,76 Mi & R\$ 721,45 Mi & R\$ 171,72 Mi & R\$ 1.358,93 Mi \\
Mar & R\$ 494,71 Mi & R\$ 458,07 Mi & R\$ 167,87 Mi & R\$ 1.120,65 Mi \\
\rowcolor{corfundo} Abr & R\$ 562,91 Mi & R\$ 525,78 Mi & R\$ 155,18 Mi & R\$ 1.243,87 Mi \\
Mai & R\$ 520,81 Mi & R\$ 638,69 Mi & R\$ 195,65 Mi & R\$ 1.355,15 Mi \\
\rowcolor{corfundo} Jun & R\$ 557,93 Mi & R\$ 640,98 Mi & R\$ 162,83 Mi & R\$ 1.361,74 Mi \\
Jul & R\$ 572,89 Mi & R\$ 419,66 Mi & R\$ 206,02 Mi & R\$ 1.198,57 Mi \\
\rowcolor{corfundo} Ago & R\$ 597,87 Mi & R\$ 543,05 Mi & R\$ 211,15 Mi & R\$ 1.352,07 Mi \\
Set & R\$ 566,48 Mi & R\$ 438,14 Mi & R\$ 189,35 Mi & R\$ 1.193,97 Mi \\
\bottomrule
\end{tabular}%
}%
}
  \par\vspace{1.5pt}{\fontsize{6.0}{7.2}\selectfont\color{cinzasub}\textit{Fonte: Boletim Tributário SEFAZ-TO e Secretaria do Tesouro Nacional (STN).}}
\end{minipage}
\par\vspace{0.08cm}
\end{minipage}\hfill
{\color{cinzaborda}\vrule width 0.4pt}\hfill
\begin{minipage}[t]{0.482\textwidth}
\par\vspace{0.08cm}
\noindent\begin{minipage}{\linewidth}
  \noindent{\fontsize{7.8}{9.5}\selectfont\bfseries\color{corprimaria}\rule{0.10cm}{0.20cm}\hspace{4pt}Arrecadação de ICMS por Segmento Econômico no Tocantins (Setembro/2026)}\par\vspace{2pt}
  {\scriptsize\setlength{\tabcolsep}{4.0pt}%
\noindent\resizebox{\linewidth}{!}{%
\begin{tabular}{lrr}
\toprule
\rowcolor{corprimaria} \textbf{\color{white}Segmento Econômico do ICMS (Painel)} & \textbf{\color{white}Arrecadação no Mês} & \textbf{\color{white}Participação no ICMS (\%)} \\
\midrule
Combustíveis & R\$ 172,81 Mi & 30,51\% \\
\rowcolor{corfundo} Comércio & R\$ 172,77 Mi & 30,50\% \\
Indústria & R\$ 103,08 Mi & 18,20\% \\
\rowcolor{corfundo} Energia Elétrica & R\$ 63,73 Mi & 11,25\% \\
Demais Atividades & R\$ 30,11 Mi & 5,32\% \\
\rowcolor{corfundo} Comunicação & R\$ 9,20 Mi & 1,62\% \\
Transportes & R\$ 8,27 Mi & 1,46\% \\
\rowcolor{corfundo} Agropecuária & R\$ 6,52 Mi & 1,15\% \\
\bottomrule
\end{tabular}%
}%
}
  \par\vspace{1.5pt}{\fontsize{6.0}{7.2}\selectfont\color{cinzasub}\textit{Fonte: SEFAZ-TO — Boletim Tributário Estadual.}}
\end{minipage}
\par\vspace{0.08cm}
\end{minipage}
\par\vspace{0.08cm}

\bannersecaoperiodico{2}{BALANÇO RETROSPECTIVO (12 MESES) E CENÁRIOS PROSPECTIVOS}{Séries Anualizadas, Consolidação Plurianual, Capilaridade Regional e Perspectivas da Equipe Técnica}
\par\vspace{0.08cm}
\noindent\begin{minipage}{\linewidth}
  \noindent{\fontsize{7.8}{9.5}\selectfont\bfseries\color{corprimaria}\rule{0.10cm}{0.20cm}\hspace{4pt}Consolidado Plurianual das Receitas Estaduais do Tocantins (2023–2026)}\par\vspace{2pt}
  {\scriptsize\setlength{\tabcolsep}{4.5pt}%
\noindent\resizebox{\linewidth}{!}{%
\begin{tabular}{lrrrrr}
\toprule
\rowcolor{corprimaria} \textbf{\color{white}Exercício} & \textbf{\color{white}Receita Total} & \textbf{\color{white}ICMS Acumulado} & \textbf{\color{white}FPE Acumulado} & \textbf{\color{white}IPVA + Outros} & \textbf{\color{white}Var. Anual ICMS} \\
\midrule
2023 (Jan-Dez) & R\$ 12.120,0 Mi & R\$ 5.028,9 Mi & R\$ 5.238,3 Mi & R\$ 2.368,4 Mi & - \\
\rowcolor{corfundo} 2024 (Jan-Dez) & R\$ 13.756,3 Mi & R\$ 5.949,4 Mi & R\$ 5.744,1 Mi & R\$ 2.623,4 Mi & +18.31\% \\
2025 (Jan-Dez) & R\$ 14.940,0 Mi & R\$ 6.407,9 Mi & R\$ 6.228,3 Mi & R\$ 2.913,3 Mi & +7.71\% \\
\rowcolor{corfundo} 2026 (Acum. Jan-Set) & R\$ 11.522,1 Mi & R\$ 4.871,0 Mi & R\$ 4.994,0 Mi & R\$ 2.023,6 Mi & +3.65\% \\
\bottomrule
\end{tabular}%
}%
}
  \par\vspace{1.5pt}{\fontsize{6.0}{7.2}\selectfont\color{cinzasub}\textit{Fonte: SEFAZ-TO e STN.}}
\end{minipage}
\par\vspace{0.08cm}

\vspace{0.03cm}
\begin{multicols}{2}
\secaoperiodico{2.1. Vetores Estruturais e Perspectivas para o Próximo Ciclo}
{\small Para o fechamento do exercício fiscal, o ingresso sazonal das receitas de fim de ano e a estabilidade da base de contribuintes ativos apontam para o cumprimento confortável das metas arrecadatórias estaduais.}

\secaoperiodico{2.2. Transmissão Intersetorial e Capilaridade Municipal}
{\small O acompanhamento sob a ótica acumulada em 12 meses evidencia que a dinâmica econômica do Tocantins apoia-se na complementaridade entre os polos regionais consolidados (Palmas, Araguaína, Gurupi, Porto Nacional, Paraíso do Tocantins e Colinas do Tocantins) e os municípios de expansão agroindustrial e logística ao longo do corredor da BR-153 e da Ferrovia Norte-Sul. A depuração de efeitos sazonais de curto prazo permite identificar ganhos persistentes de escala produtiva, formalização empresarial e densidade salarial.}

\secaoperiodico{2.3. Síntese Institucional e Metodológica (Ótica de 12 Meses)}
{\small A leitura conjunta das séries acumuladas em quatro trimestres e das comparações interanuais confirma a solidez dos fundamentos regionais do Tocantins neste eixo temático. O monitoramento contínuo das bases oficiais pelo Jornal Econômico do Tocantins — integrando extração automatizada via APIs públicas, tratamento de microdados e revisão analítica da equipe — permite distinguir oscilações transitórias de mudanças estruturais de tendência, subsidiando decisões de investimento privado e formulação de políticas públicas estaduais e municipais.}
\end{multicols}


\vspace{0.12cm}

% ==============================================================================
% EXPEDIENTE INSTITUCIONAL E NOTA METODOLÓGICA
% ==============================================================================
\noindent\begin{tikzpicture}
  \node[
    fill=corfundo,
    draw=cinzaborda,
    line width=0.5pt,
    rounded corners=2pt,
    inner sep=6pt,
    text width=\textwidth-13pt,
    font=\scriptsize\color{cinzasub}
  ] {
    \textbf{\color{corprimaria}Expediente e Nota Metodológica:} Consolida os indicadores fiscais do painel: Receita Tributária Própria (ICMS, IPVA, ITCD, IRRF e Taxas) extraída dos Boletins Tributários da SEFAZ-TO e repasses constitucionais do FPE (STN).\par\vspace{2pt}
    \textbf{\color{corprimaria}Equipe Responsável:} Lorenzo Costa | Graduado em Economia — Jornal Econômico do Tocantins \enspace$\vert$\enspace \textbf{\color{corprimaria}Publicação Institucional:} Núcleo de Conjuntura Econômica — \textit{Jornal Econômico do Tocantins}.
  };
\end{tikzpicture}

\end{document}
